\element{1.3.1}{
  \begin{questionmult}{OscillateurWienComposition}\bareme{haut=2,p=0}
    De quoi est constitué un oscillateur de Wien ?
    \begin{multicols}{2}
      \begin{reponses}
        \bonne{Un amplificateur}
        \bonne{Un filtre passe bande}
        \mauvaise{Un comparateur à hystérésis}
        \mauvaise{Un intégrateur}
      \end{reponses}
    \end{multicols}
  \end{questionmult}
}
\element{1.3.1}{
  \begin{questionmult}{OscillateurRelaxationComposition}\bareme{haut=2,p=0}
    De quoi est constitué un oscillateur à relaxation ?
    \begin{multicols}{2}
      \begin{reponses}
        \mauvaise{Un amplificateur}
        \mauvaise{Un filtre passe bande}
        \bonne{Un comparateur à hystérésis}
        \bonne{Un intégrateur}
      \end{reponses}
    \end{multicols}
  \end{questionmult}
}
\setgroupmode{1.3.1}{cyclic}

\element{1.3.2}{
  \begin{questionmult}{WienDemarrage}\bareme{haut=2,p=0}
    Pour que des oscillations démarrent dans un oscillateur quasi-sinusoïdal, il faut que
    \begin{multicols}{2}
      \begin{reponses}
        \mauvaise{L'ALI soit dans un montage instable}
        \mauvaise{Le produit des fonctions de transfert de chaque système soit égal à $1$}
        \mauvaise{Chaque élément du système bouclé soient instable}
        \bonne{Le système bouclé soit instable}
      \end{reponses}
    \end{multicols}
  \end{questionmult}
}
\element{1.3.2}{
  \begin{questionmult}{WienExistance}\bareme{haut=2,p=0}
    Pour que les oscillations sinusoïdales existent dans un oscillateur quasi-sinusoïdal, il faut que
    \begin{multicols}{2}
      \begin{reponses}
        \mauvaise{L'ALI soit dans un montage instable}
        \bonne{Le produit des fonctions de transfert de chaque système soit égal à $1$}
        \mauvaise{Chaque élément du système bouclé soient instable}
        \mauvaise{Le système bouclé soit instable}
      \end{reponses}
    \end{multicols}
  \end{questionmult}
}
\setgroupmode{1.3.2}{cyclic}